\chapter{Espacios de Hilbert}

\begin{definicion}
    Sea $X$ un espacio vectorial sobre $\mathbb{R}$. Se dice que $X$ es un \textbf{espacio con producto interno} si existe una aplicación:
    \Func{(\cdot, \cdot)}{X \times X}{\mathbb{R}}{(x,y)}{x\cdot y=(x,y)}
    tal que:
    \begin{itemize}
        \item[i)] $(x + y, z) = (x, z) + (y, z) \quad \forall x, y, z \in X$
        \item[ii)] $(\alpha x, y) = \alpha(x, y) \quad \forall x, y \in X, \ \forall \alpha \in \mathbb{R}$
        \item[iii)] $(x, y) = (y, x) \quad \forall x, y \in X$
        \item[iv)] $(x, x) \geq0$ \ y \ $(x, x) = 0 \iff x = 0 \quad \forall x \in X$
    \end{itemize}
    Sobre $\mathbb{C}$, la condición $iii)$ se sustituye por $iii)'$: $(x, y) = \overline{(y, x)}$ y por lo tanto:
        $$(x, \alpha y + \beta z) = \bar{\alpha}(x, y) + \bar{\beta}(x, z) \quad \forall \alpha, \beta \in \mathbb{C} \text{ \textit{(producto sesquilineal)}}$$
    Destacar que la diferencia con el producto interno en $\R$ es que esta vez el producto interno es lineal en la primera variable y conjugado lineal en la segunda.
\end{definicion}

\observacion{
    Las condiciones $i)$ y $ii)$ nos dicen que $(\cdot, \cdot)$ es lineal en la primera variable y, junto con $iii)$, se tiene que $(\cdot, \cdot)$ es bilineal sobre $\mathbb{R}$.
}

\begin{observacion}
    $(X, (\cdot, \cdot))$ es un espacio normado con norma $\|x\| := (x, x)^{1/2}$. \\
\end{observacion}

\begin{prop}[Desigualdad de Cauchy-Schwarz]
    Sea $X$ un espacio con producto interno, entonces se cumple:
    \[ |(x, y)| \le \|x\| \cdot \|y\| \quad \forall x, y \in X \]
\end{prop}

\begin{definicion}
    Un espacio $(X, (\cdot, \cdot))$ con producto interno se dice \textbf{espacio de Hilbert} \mbox{$\mathcal{H} = (X, (\cdot, \cdot))$} si es completo.
\end{definicion}


\begin{ejemplo}
    \begin{itemize}
        \item $\ell_2$ es un espacio de Hilbert: $(x, y) := \sum_{k=0}^{+\infty} x_k \cdot \overline{y_k}$
        \item $L^2(A)$ es un espacio de Hilbert: $(f, g) := \int_A f \bar{g} \, d\mu$
    \end{itemize}
\end{ejemplo}

\noindent
Buscamos un criterio que nos permita entender cuándo un espacio normado tiene una norma que deriva de un producto interno.

\begin{teo}[Ley del paralelogramo]\label{teo:paralelogramo}
    Sea $(X, \|\cdot\|)$ un espacio normado, entonces
    \begin{equation*}
        X \text{ espacio con producto interno con } \|x\|^2 := (x, x) \iff \|x+y\|^2 + \|x-y\|^2 = 2\|x\|^2 + 2\|y\|^2 \quad \forall x, y \in X
    \end{equation*}

\begin{proof}
    \begin{description}
        \item[$\Longrightarrow)$] Sean $x, y \in X$, entonces:
        \begin{equation*}
        \begin{cases}
            \|x+y\|^2 = (x+y, x+y) = \|x\|^2 + 2(x, y) + \|y\|^2 \\
            \|x-y\|^2 = (x-y, x-y) = \|x\|^2 - 2(x, y) + \|y\|^2
            \end{cases}
            \implies \|x+y\|^2 + \|x-y\|^2 = 2(\|x\|^2 + \|y\|^2)
        \end{equation*}
        \item[$\Longleftarrow)$] Supongamos que vale la ley del paralelogramo, definimos un producto interno
        \begin{equation*}
            (x, y) := \frac{1}{4} (\|x+y\|^2 - \|x-y\|^2) \quad \forall x, y \in X
        \end{equation*}
        demostremos que $(\cdot, \cdot)$ es un producto interno
        \begin{itemize}
            \item \textbf{Simétrico:} $(y, x) = (x, y)$
            \item \textbf{Norma:} $\|x\| = (x, x)^{1/2} = \left( \frac{1}{4} \cdot 4\|x\|^2 \right)^{1/2} = \|x\| \ge 0$
            \item \textbf{Linealidad en la primera componente:} 
            Queremos demostrar primero la aditividad: $(x+z, y) = (x, y) + (z, y)$.
            Aplicando la definición del producto interno candidato, tenemos:
            \begin{align*}
                (x, y) + (z, y) &= \frac{1}{4} \left( \|x+y\|^2 - \|x-y\|^2 \right) + \frac{1}{4} \left( \|z+y\|^2 - \|z-y\|^2 \right) \\
                &= \frac{1}{4} \left( \left[ \|x+y\|^2 + \|z+y\|^2 \right] - \left[ \|x-y\|^2 + \|z-y\|^2 \right] \right) \quad (*)
            \end{align*}
            Por la ley del paralelogramo, sabemos que para cualesquiera dos vectores $a, b \in X$ se cumple que $\|a+b\|^2 + \|a-b\|^2 = 2\|a\|^2 + 2\|b\|^2$, lo que equivale a:
            \begin{equation*}
                \|a\|^2 + \|b\|^2 = \frac{1}{2} \left( \|a+b\|^2 + \|a-b\|^2 \right)
            \end{equation*}
            Aplicamos esta propiedad a los dos bloques de el término $(*)$:
            \begin{itemize}
                \item Para el primer bloque, tomamos $a = x+y$ y $b = z+y$:
                \begin{equation*}
                    \|x+y\|^2 + \|z+y\|^2 = \frac{1}{2} \left( \|x+z+2y\|^2 + \|x-z\|^2 \right)
                \end{equation*}
                \item Para el segundo bloque, tomamos $a = x-y$ y $b = z-y$:
                \begin{equation*}
                    \|x-y\|^2 + \|z-y\|^2 = \frac{1}{2} \left( \|x+z-2y\|^2 + \|x-z\|^2 \right)
                \end{equation*}
            \end{itemize}
            Sustituyendo ambos bloques de nuevo en $(*)$, el término común $\|x-z\|^2$ se cancela al restar, quedando:
            \begin{align*}
                (x, y) + (z, y) &= \frac{1}{4} \cdot \frac{1}{2} \left( \|x+z+2y\|^2 - \|x+z-2y\|^2 \right) \\
                &= \frac{1}{8} \left( \|x+z+2y\|^2 - \|x+z-2y\|^2 \right) \quad (**)
            \end{align*}
            Por otro lado, si aplicamos la ley del paralelogramo directamente a los vectores $a = x+z$ y $b = 2y$, obtenemos:
            \begin{equation*}
                \|x+z+2y\|^2 + \|x+z-2y\|^2 = 2\|x+z\|^2 + 2\|2y\|^2 = 2\|x+z\|^2 + 8\|y\|^2
            \end{equation*}
            Y si evaluamos $(x+z, 2y)$ usando que por simetría y aditividad preliminar (haciendo $x=z$) se cumple que $(2x, y) = 2(x,y)$, la relación $(**)$ se reduce exactamente a:
            \begin{equation*}
                \frac{1}{8} \left( \|x+z+2y\|^2 - \|x+z-2y\|^2 \right) = 2 \left( \frac{1}{4} \left( \|x+z+y\|^2 - \|x+z-y\|^2 \right) \right) = (x+z, y)
            \end{equation*}
            Por lo tanto, queda demostrada la aditividad: $(x+z, y) = (x, y) + (z, y)$.

            \vspace{0.2cm}
            \item \textbf{Homogeneidad:} Queremos demostrar que $(\alpha x, y) = \alpha (x, y)$ para todo $\alpha \in \mathbb{R}$.
            \begin{itemize}
                \item Por la aditividad que acabamos de probar, mediante inducción matemática es inmediato que $(n x, y) = n(x, y)$ para todo $n \in \mathbb{N}$, y se extiende fácilmente a los enteros $\mathbb{Z}$ usando que $(-x, y) = -(x,y)$ por la simetría de la norma.
                \item Para los números racionales $q = \frac{m}{n} \in \mathbb{Q}$, tenemos:
                \begin{equation*}
                    n \left( \frac{m}{n} x, y \right) = \left( m x, y \right) = m (x, y) \implies \left( \frac{m}{n} x, y \right) = \frac{m}{n} (x, y)
                \end{equation*}
                \item Para un número real cualquiera $\alpha \in \mathbb{R}$, tomamos una sucesión de racionales $(q_k) \subseteq \mathbb{Q}$ tal que $q_k \to \alpha$. Como la norma es una aplicación continua, el producto interno candidato $(\cdot, \cdot)$ también es continuo respecto a sus variables. Por lo tanto, pasando al límite:
                \begin{equation*}
                    (\alpha x, y) = \lim_{k \to \infty} (q_k x, y) = \lim_{k \to \infty} q_k (x, y) = \alpha (x, y)
                \end{equation*}
            \end{itemize}
            Habiendo demostrado la simetría, la definición compatible de la norma y la linealidad completa, concluimos que $(\cdot, \cdot)$ es un producto interno genuino en $X$.
        \end{itemize}
    \end{description}
\end{proof}
\end{teo}

\begin{ejercicio}
    $\ell_1$ no tiene producto interno. \\

    \noindent
    Sean $x = (x_m), y = (y_m) \in \ell_1, y \neq 0$, y donde $\exists k_0 \in \mathbb{N}$ tal que $x_{k_0}, y_{k_0}$ tienen símbolos opuestos, entonces
    \begin{equation*}
        \|x+y\| = \sum_{k=1}^{\infty} |x_k + y_k| < \sum_{k=1}^{\infty} |x_k| + \sum_{k=1}^{\infty} |y_k| = \|x\| + \|y\|
    \end{equation*}
    donde la desigualdad estricta se debe a que $x_{k_0}$ y $y_{k_0}$ tienen símbolos opuestos, de lo cual 
    \begin{equation*}
        \|x+y\| + \|x-y\| < 2\|x\| + 2\|y\|
    \end{equation*}
    \textit{(Nota: Al no cumplirse la igualdad de la ley del paralelogramo, la norma no proviene de un producto interno).}
\end{ejercicio}

\begin{ejercicio}
    $(C[0,\pi], \|\cdot\|_{\infty})$ no tiene producto interno. \\
\end{ejercicio}

\begin{prop}
    $z = 0 \iff (x, z) = 0 \quad \forall x \in X$\\

\begin{proof}
    \begin{description}
        \item[$\Longrightarrow)$] $(x, z) =(x,0)= (x,x-x) = (x, x) - (x, x) = 0 \quad \forall x \in X$
        \item[$\Longleftarrow)$] Si $(x, z) = 0$ para todo $x \in X$, entonces $(z, z) = 0$ y por lo tanto $z = 0$.
    \end{description}
\end{proof}
\end{prop}

\begin{ejercicio}
    Sea $\{z_m\} \in X$ tal que $(z_m, x) \longrightarrow (z, x)$ y $\|z_m\| \longrightarrow \|z\|$, entonces $z_m \to z$ (fuerte). \\

    \noindent
    Por hipótesis $(z_m, x) \longrightarrow (z, x) \quad \forall x \in X \iff (z_m - z, x) \longrightarrow 0 \quad \forall x \in X$, entonces tomando $x=z$, tenemos que
    $$(z_m, z) \to (z, z) = \|z\|^2 \text{ y } \|z_m\|^2 \to \|z\|^2$$
    Además con el siguiente desarrollo
    $$\|z_m - z\|^2 = (z_m - z, z_m - z) = \|z_m\|^2 - (z_m, z) - (z, z_m) + \|z\|^2$$
    podemos concluir que 
    $$\lim \|z_m - z\|^2 = \|z\|^2 - \|z\|^2 - \|z\|^2 + \|z\|^2 = 0 \implies z_m \to z$$
\end{ejercicio}

\observacion{
    Sea $\bar{x}\in X$ tal que $F_{\bar{x}}(y) = (\bar{x}, y) \quad \forall y \in Y$, entonces $F_{\bar{x}}$ es lineal y continuo, veámoslo
    \begin{itemize}
        \item La linealidad se hereda de la linealidad del producto interno.
        \item $|F_{\bar{x}}(y)| = |(\bar{x}, y)| \le \|\bar{x}\| \cdot \|y\| \implies \|F_{\bar{x}}\| \le \|\bar{x}\|$, y entonces tomando $y = \bar{x}$, tenemos que
        $$F_{\bar{x}}(y) = F_{\bar{x}}(\bar{x}) = (\bar{x}, \bar{x}) = \|\bar{x}\| \cdot \|\bar{x}\| \implies \|F_{\bar{x}}\| = \|\bar{x}\|$$
    \end{itemize}
}


\begin{teo}
    Todo espacio con producto interno $(X, (\cdot, \cdot))$ es uniformemente convexo.
\begin{proof}
    Sean $x, y \in X$ tales que $\|x\| \le 1$, $\|y\| \le 1$ y $\|x - y\| > \varepsilon$. Demostremos que $\frac{1}{2} \|x + y\| \le 1 - \delta$ para algún $\delta = \delta_\varepsilon$.

    $$\frac{1}{4} \|x + y\|^2 = \frac{2\|x\|^2 + 2\|y\|^2}{4} - \frac{\|x - y\|^2}{4} \le \frac{1}{2} + \frac{1}{2} - \frac{1}{4} \varepsilon^2 = 1 - \frac{\varepsilon^2}{4} < 1$$

    Sea $\delta := \frac{\varepsilon^2}{4} \implies \|\frac{1}{2}(x + y)\|^2 = 1 - \delta, \quad \delta > 0$.
\end{proof}
\end{teo}

\begin{ejercicio}
    Sea $C$ un conjunto cerrado y convexo de $\mathcal{H}$ espacio de Hilbert, entonces demostrar que $C$ admite un único punto de mínima distancia. \\ \\
    
    \noindent
    Sea $\mathcal{H}$ un espacio de Hilbert $\implies \mathcal{H}$ es uniformemente convexo y, por lo tanto, $\mathcal{H}$ es un espacio reflexivo, de lo cual se sigue que existe un punto de norma mínima. Sean $x_0, x_1$ puntos de norma mínima, es decir, $\|x_0\| = \|x_1\| = d$. Supongamos por el contrario (P.A.) que $x_0 \neq x_1$. \\ \\
    \noindent
    Por la convexidad de $C$, el punto medio $\frac{x_0 + x_1}{2} \in C$, por lo tanto su norma debe ser mayor o igual al ínfimo:
    \[ \left\| \frac{x_0 + x_1}{2} \right\| \ge d \]
    Por otro lado, utilizando la \textbf{ley del paralelogramo}:
    \[ \left\| \frac{x_0 + x_1}{2} \right\|^2 = \frac{1}{4} \|x_0 + x_1\|^2 = \frac{1}{4} \left( 2\|x_0\|^2 + 2\|x_1\|^2 - \|x_0 - x_1\|^2 \right) = \]
    \[ = \frac{1}{2} \|x_0\|^2 + \frac{1}{2} \|x_1\|^2 - \frac{1}{4} \|x_0 - x_1\|^2 = d^2 - \frac{1}{4} \|x_0 - x_1\|^2 < d^2 \]
    puesto que hemos supuesto $x_0 \neq x_1 \implies \|x_0 - x_1\| > 0$. \\
    \noindent
    Esto implica que $\left\| \frac{x_0 + x_1}{2} \right\| < d$, lo cual es un **ABSURDO**, de donde se concluye que $x_0 = x_1$. \\
\end{ejercicio}

\begin{teo}[Teorema de proyección]
    Sea $\mathcal{H}$ un espacio de Hilbert, y sea $C \subseteq \mathcal{H}$ un subconjunto no vacío, cerrado y convexo, entonces para cada $x \in \mathcal{H}$ existe un único $x_0 \in C$ tal que:
    \[ \|x - x_0\| = \inf_{y \in C} \|x - y\| \]
    Además, $x_0$ está caracterizado por la siguiente propiedad
    \begin{equation*}
        x_0 \text{ es el único punto de } C \text{ tal que } (x - x_0, y - x_0) \le 0 \quad \forall y \in C
    \end{equation*}
    El punto $x_0 \in X$ se denomina proyección de $x$ sobre $C$ y se indica como $P_C(x)$.
\begin{proof}
    $\mathcal{H}$ es espacio de Hilbert $\implies \mathcal{H}$ es reflexivo y $C$ es cerrado y convexo, entonces $\exists! x_0 \in C$ que minimiza la distancia $\|x - y\|$. Mostremos la caracterización de $x_0$.
    \begin{description}
        \item[$\Longrightarrow)$] Supongamos $x_0 \in C : \|x - x_0\| = \inf\limits_{y \in C} \|x - y\|$ y sea $z \in \{(1-t)x_0 + ty \mid t \in [0,1]\} \subseteq C$, entonces:
        \[ \|x - x_0\| \le \|x - z\| = \|x - x_0 - t(y - x_0)\| \iff \|x - x_0\|^2 \le \|x - x_0 - t(y - x_0)\|^2 \]
        de lo cual:
        \[ \|x - x_0\|^2 \le \|x - x_0\|^2 + t^2 \|y - x_0\|^2 - 2t(x - x_0, y - x_0) \iff t \|y - x_0\|^2 - 2(x - x_0, y - x_0) \ge 0 \quad \forall t \in \mathbb{R} \]
        haciendo $t \to 0$:
        \[ (x - x_0, y - x_0) \le 0 \]
        \item[$\Longleftarrow)$] Sea $x_0 \in C$ el único punto tal que $(x - x_0, y - x_0) \le 0 \quad \forall y \in C$, entonces:
        \[ \|x - x_0\|^2 = \|x - y + y - x_0\|^2 = \|x - y\|^2 + \|y - x_0\|^2 + 2(x - y, y - x_0) \le \]
        \[ \le \|x - y\|^2 + 2 \underbrace{(x - x_0, y - x_0)}_{\le 0} - \|y - x_0\|^2 \le \|x - y\|^2 \quad \forall y \in C \]
    \end{description}
\end{proof}
\end{teo}

\begin{ejercicio}
    La proyección $P_C$ es Lipschitz de constante $L=1$, es decir
    \[ \|P_C(z_1) - P_C(z_2)\| \le \|z_1 - z_2\| \quad \forall z_1, z_2 \in \mathcal{H} \]
\end{ejercicio}

\section{Ortogonalidad en espacios de Hilbert}

\begin{definicion}
    Sean $x, y \in \mathcal{H}$, se dice que $x$ es \textbf{ortogonal} a $y$ si $(x, y) = 0$, y escribiremos $x \perp y$.
\end{definicion}

\begin{prop}
    Si $x \perp y$, entonces se tiene 
    \begin{itemize}
        \item[i)] $y \perp x$
        \item[ii)] $x \perp (-y)$
        \item[iii)] $x \perp x \iff x = 0$
    \end{itemize}
\end{prop}

\begin{definicion}
    Sean $A, B \subseteq \mathcal{H}$ subconjuntos. Se dice que $x \in \mathcal{H}$ es \textbf{ortogonal a $A$} ($x \perp A$) si 
    $$(x, a) = 0 \quad \forall a \in A$$ 
    \noindent
    Se dice que $A$ es \textbf{ortogonal a $B$} ($A \perp B$) si 
    $$(a, b) = 0 \quad \forall a \in A, \ \forall b \in B$$ 
    \noindent
    El \textbf{conjunto ortogonal a $A$} se define como:
    \[ A^\perp := \{ x \in \mathcal{H} \mid (x, a) = 0 \quad \forall a \in A \} \]
\end{definicion}

\begin{prop}\label{prop:ortogonalidad}
    Se cumplen las siguientes propiedades:
    \begin{enumerate}
        \item $0^{\perp} = \mathcal{H}$ \ y \ $\mathcal{H}^{\perp} = \{0\}$
        \item $A \cap A^{\perp} \subseteq \{0\}$
        \item (\textit{Teorema de Pitágoras}) si $x \perp y$ entonces $\|x + y\|^2 = \|x\|^2 + \|y\|^2$
        \item Sean $A, B$ subconjuntos con $A \subseteq B$ entonces $B^{\perp}\subseteq A^{\perp}$
    \end{enumerate}
\end{prop}

\begin{ejercicio}
    Sea $A \subseteq \mathcal{H}$ un subconjunto, entonces $A^{\perp}$ es un subespacio cerrado de $\mathcal{H}$. \\
\end{ejercicio}

\begin{prop}
    Sea $x \in \mathcal{H}$ y sea $x_0 := P_{M}(x)$, donde $M\subseteq\mathcal{H}$ no vacío, cerrado y convexo, entonces $x_0$ está caracterizado por estas dos condiciones
    \begin{itemize}
        \item[i)] $x_0 \in M$
        \item[ii)] $(x - x_0, y) = 0 \quad \forall y \in M \quad (x - x_0 \in M^\perp)$.
    \end{itemize}
    es decir, $x_0=P_M(x)$ si y solo si se cumplen $i)$ y $ii)$. \\
    \noindent
    Además, la función $P_M$
    \Func{P_M}{\mathcal{H}}{M}{x}{P_M(x)}
    es un operador lineal y continuo, y se define como \textbf{proyección ortogonal sobre $M$}.
\begin{proof}
    \begin{description}
        \item[$\Longrightarrow)$] $x_0 = P_M(x)$ y por lo tanto $(x - x_0, y - x_0) \le 0 \quad \forall y \in M$. Al ser $M$ un subespacio, entonces $y + x_0 \in M$ y por lo tanto se tiene:
        \[ (x - x_0, y + x_0 - x_0) \le 0 \implies (x - x_0, y) \le 0 \quad \forall y \in M \]
        de lo cual:
        \[ (x - x_0, -y) \le 0 \iff -(x - x_0, y) \le 0 \iff (x - x_0, y) \ge 0 \]
        \[ \implies (x - x_0, y) = 0 \quad \forall y \in M \]
        \item[$\Longleftarrow)$] Supongamos ciertas i), ii), entonces
        $$(x - x_0, y) = 0 \quad \forall y \in M \text{ y } y - x_0 \in M \implies (x - x_0, y - x_0) = 0 \quad \forall y \in M$$
        donde en la condición $y - x_0 \in M$ se ha utilizado que $x_0\in M$. Teniendo el resultado anterior, se concluye que $x_0 = P_M(x)$, al usar la caracterización del Teorema de proyección. \\
    \end{description}
    Demostremos ahora que $P_M$ es un operador lineal y continuo:
    \begin{itemize}
        \item \tbf{Linealidad:} Sean $x_1, x_2 \in \mathcal{H}$, entonces \mbox{$(x_1 - P_M(x_1), y) = 0$ y $(x_2 - P_M(x_2), y) = 0 \quad \forall y \in M$}. Sumando:
        \[ (x_1 + x_2 - [P_M(x_1) + P_M(x_2)], y) = 0 \quad \forall y \in M \overset{unicidad}{\implies} P_M(x_1 + x_2) = P_M(x_1) + P_M(x_2) \]
        Analógamente:
        \begin{align*}
            (x_1 - P_M(x_1), y) = 0 &\implies \alpha(x_1 - P_M(x_1), y) = 0 \\
            &\implies (\alpha x_1 - \alpha P_M(x_1), y) = 0 \\
            &\implies \alpha P_M(x_1) = P_M(\alpha x_1)
        \end{align*}
        \item \tbf{Continuidad:} se deduce de la propiedad de Lipschitz (con $L=1$).
    \end{itemize}
\end{proof}
\end{prop}

\begin{definicion}
    Sea $X$ un espacio normado, $X$ es la \textbf{suma directa} de dos subespacios $A, B$ si
    \begin{itemize}
        \item $A, B$ son cerrados.
        \item $A \cap B = \{0\}$
        \item $A + B = X$
    \end{itemize}
    y diremos $X = A \oplus B$.
\end{definicion}

\begin{teo}[Teorema de descomposición ortogonal]
    Sea $M$ un subespacio cerrado, entonces 
    $$\mathcal{H} = M \oplus M^\perp$$
\begin{proof}
    $M^\perp$ es subespacio cerrado, por lo tanto sabemos que $0\in M, M^\perp$, por lo que $M \cap M^\perp = \{0\}$, usando la Proposición \ref{prop:ortogonalidad}. Mostremos que $M + M^\perp = \mathcal{H}$, veámoslo. \\

    \noindent
    Sea $x \in \mathcal{H}$, entonces $P_M(x) =: x_0 \in \mathcal{H}$ y por lo tanto $x = x_0 + (x - x_0)$. 
    Como \mbox{$(x - x_0, y) = 0 \quad \forall y \in M$}, entonces $x - x_0 \in M^\perp$, y siendo $x_0 \in M$, se tiene:
    \[ x = \underbrace{x_0}_{\in M} + \underbrace{(x - x_0)}_{\in M^\perp}\]
    Destacar que hemos usado la caracterización del Teorema de proyección donde la característica de $M$ de ser convexo viene implícita al ser un subespacio. 
\end{proof}
\end{teo}

\observacion{
    Si $\mathcal{H} = M \oplus M^\perp$, la descomposición es única, veámoslo. Si $x = \overbrace{a}^{\in M} + \overbrace{b}^{\in M^\perp}$, entonces 
    $$x - a = b \in M^\perp \implies (x - a, y) = 0 \quad \forall y \in M \implies a = P_M(x) = x_0 \text{ y } b = x - x_0$$
    tenemos por tanto que $a=x_0$ es único y $b = x - x_0$ también es único, por la noción de que la proyección de $x$, es decir, $x_0$, es única.
}

\observacion{
    Sea $A \subseteq \mathcal{H}$ un subespacio, entonces se cumple:
    \begin{itemize}
        \item[i)] $A \subseteq (A^\perp)^\perp$
        \item[ii)] $(A^\perp)^\perp = \overline{A}$. En particular si $A$ es cerrado, entonces $A = (A^\perp)^\perp$
    \end{itemize}
}

\begin{ejercicio}
    Sea $M \subseteq \mathcal{H}$ un subespacio cerrado, entonces:
    \begin{itemize}
        \item[i)] $(P_M(x), y) = (x, P_M(y)) \quad \forall x, y \in \mathcal{H}$
        \item[ii)] $P_M^2(x) = P_M(x) \quad \forall x \in \mathcal{H}$
    \end{itemize}
\end{ejercicio}

\observacion{
    Sea $X$ un espacio con producto interno y sea 
    \Func{T}{X}{X^*}{y}{F_y}
    donde $F_y(x) = (x, y) \quad \forall x \in X$. Entonces:

    \begin{enumerate}
        \item \textbf{$T$ es inyectiva:} si $y_1 \neq y_2 \implies \exists x_0 \in X : (x_0, y_1) \neq (x_0, y_2)$. 
        De lo contrario, si no existiera tal $x_0 \in X$:
        \[ (x, y_1) = (x, y_2) \quad \forall x \in X \iff (x, y_1 - y_2) = 0 \quad \forall x \in X  \]
        entonces tomando $x=y_1 - y_2\in X$ se concluye que 
        $$(y_1 - y_2, y_1 - y_2) = 0 \implies y_1 - y_2 = 0 \implies y_1 = y_2$$
        lo cual es un \textbf{ABSURDO}.
        
        \item \textbf{$T$ es isometría:} $\|F_y\| = \|y\|$, como ya se ha visto anteriormente, y como una isometría conserva las distancias, es decir, las normas, entonces $T$ es isometría.
    \end{enumerate}
}

\begin{teo}[Teorema de Riesz]
    Sea $\mathcal{H}$ un espacio de Hilbert, entonces 
    \begin{equation*}
        \forall F \in \mathcal{H}^* \quad \exists! y_F \in \mathcal{H} \text{ tal que } F(x) = (x, y_F) \quad \forall x \in \mathcal{H}
    \end{equation*}
    es decir, $T$ es sobreyectiva.

\begin{proof}
    Busquemos ese $y_F$ de ese tipo. Sea 
    $$N_F := \text{Ker}(F) = \{x \in \mathcal{H} \mid F(x) = 0\}$$ 
    es un subespacio cerrado (por ser $F$ funcional continuo y $Ker(F)=F^{-1}(\{0\}))$ y propio de $\mathcal{H}$ (de lo contrario, si $F \equiv 0$, entonces $y_F = 0$), por lo tanto
    $$\mathcal{H} = N_F \oplus N_F^\perp$$
    Sea $x_0 \in N_F^\perp \setminus \{0\}$ y sea $x \in \mathcal{H}$. Consideremos $x - \dfrac{F(x)}{F(x_0)}x_0 \in N_F$, donde es clave que $x_0 \in N_F^\perp$ y $F(x_0) \neq 0$ (por ser $x_0 \neq 0$ y $x_0 \in N_F^\perp$), entonces
    \[ F\left( x - \frac{F(x)}{F(x_0)}x_0 \right) = F(x) - \frac{F(x)}{F(x_0)} F(x_0) = 0 \]
    por lo tanto
    \begin{equation*}
        \left( \underbrace{\left( x - \frac{F(x)}{F(x_0)}x_0 \right)}_{\in N_F}, \underbrace{x_0}_{\in N_F^\perp} \right)\overset{(*)}{=} 0\implies (x, x_0) - \left( \frac{F(x)}{F(x_0)}x_0, x_0 \right) = 0 \implies(x, x_0) = \frac{F(x)}{F(x_0)} \|x_0\|^2 \quad \forall x \in \mathcal{H} 
    \end{equation*}
    donde en $(*)$ hemos usado que son perpendiculares al pertenecer a $N_F, N_F^\perp$. Tenemos a continuación
    \[ F(x) = \frac{F(x_0)}{\|x_0\|^2} (x, x_0) \]
    y definiendo $y_F := \dfrac{\overline{F(x_0)}}{\|x_0\|^2} x_0$, tenemos que $y_F$ representa a $F$ y es único porque $T$ es inyectiva.
\end{proof}
\end{teo}

\observacion{
    Hemos visto que $T: \mathcal{H} \longrightarrow \mathcal{H}^*$ es un isomorfismo algebraico y topológico (dijimos $T$ isometría, que implica continuidad y de su inversa, y biyectiva), entonces $\mathcal{H}^*$ es un espacio de Hilbert, veámoslo $\forall F, G \in \mathcal{H}^*$
    \[ (F, G) := (y_F, y_G) = (T^{-1}(F), T^{-1}(G)) \]
    y $\|\cdot\|_{\mathcal{H}^*}$ es la misma debido al producto escalar.
    Por lo tanto, se tiene: $\mathcal{H} \xrightarrow{T_{\mathcal{H}}} \mathcal{H}^* \xrightarrow{T_{\mathcal{H}^*}} \mathcal{H}^{**}$ con $T_{\mathcal{H}}, T_{\mathcal{H}^*}$ isomorfismos $\implies T_{\mathcal{H}^*} \circ T_{\mathcal{H}} : \mathcal{H} \xrightarrow{\sim} \mathcal{H}^{**}$ es un isomorfismo y tal que $T_{\mathcal{H}^*} \circ T_{\mathcal{H}} = J : \mathcal{H} \xrightarrow{\sim} \mathcal{H}^{**}$. Se sigue que $\mathcal{H}$ es reflexivo.
}

\begin{ejemplo}
    $H_0^1(\Omega) \hookrightarrow L^2(\Omega) \simeq L^2(\Omega)^* \hookrightarrow H_0^1(\Omega)^*$ \quad \textit{atención al identificar}
\end{ejemplo}

\section{Operador adjunto}

\noindent
Sea $\mathcal{H}$ un espacio de Hilbert y sea $x \in \mathcal{H}$ fijo. Definimos 
\Func{G_x}{\mathcal{H}}{\mathbb{R} \text{ o } \mathbb{C}}{y}{(x, T(y))}
con $T$ operador lineal en $\mathcal{H}$. \\
\noindent
Como $G_x$ es un funcional lineal y, por el teorema de Riesz, $\exists! z_x \in \mathcal{H}$ tal que $G_x(y) = (z_x, y)$, es decir
\[ (x, T(y)) = (z_x, y) \]
El análisis funcional no se queda en un solo punto, sino que mira el ``comportamiento global''. Si ahora cambias el vector $x$ por otro vector $x'$, el Teorema de Riesz te dará un nuevo vector único $z_{x'}$. Esto significa que hemos creado, sin darnos cuenta, una regla de correspondencia:

\begin{itemize}
    \item A cada $x \in \mathcal{H}$ le asignamos un vector $z_x \in \mathcal{H}$ tal que $(x, T(y)) = (z_x, y) \quad \forall y \in \mathcal{H}$
    \item A esa función (ese mapeo) es a lo que bautizamos como \underline{operador adjunto} y lo denotamos como $T^*$.  
\end{itemize}
\noindent
Por lo tanto, escribimos $T^*(x)$ para indicar que ese vector es el resultado de aplicar el operador $T^*$ al vector $x$.

\begin{definicion}
    Sea $T \in BL(\mathcal{H})$, se define como \textbf{operador adjunto} de $T$ a la aplicación
    \Func{T^*}{\mathcal{H}}{\mathcal{H}}{x}{T^*(x)}
    tal que $T^*(x)$ es el único vector que representa a $G_x$, es decir
    \[ G_x(y)=(x, T(y)) = (T^*(x), y) \quad \forall x, y \in \mathcal{H} \]
\end{definicion}

\begin{prop}
    Sea $T\in BL(\mathcal{H})$, entonces 
    \begin{itemize}
        \item[i)] $T^{**} = T$
        \item[ii)] $\|T^*\| = \|T\|$, y por lo tanto $T^* \in BL(\mathcal{H})$
    \end{itemize}
\begin{proof}
    Probemos una inclusión de $ii)$, primero.
    \begin{equation*}
        \|T^*(x)\|^2 = (T^*(x), T^*(x)) = (x, T(T^*(x))) \le \|x\| \cdot \|T(T^*(x))\| \le \|x\| \cdot \|T\| \cdot \|T^*(x)\|
    \end{equation*}
    de lo cual:
    \[ \|T^*(x)\| \le \|x\| \cdot \|T\| \implies \|T^*\| \le \|T\| \]
    por lo que tenemos que $T^*$ está acotada, para ver que $T^*\in BL(\mathcal{H})$ necesitamos la linealidad. La linealidad de $T^*$ viene heredada de la estructura del producto interno y la linealidad de $T$, por lo que $T^*\in BL(\mathcal{H})$. \\
    \noindent
    Por otro lado, veámos ahora $i)$, pues sabiendo lo anterior, aplicando aqui $(**)$ otra vez la definición de operador adjunto que
    $$(x, T(y)) = (T^*(x), y) \overset{(**)}{=} (x, T^{**}(y)) \quad \forall x, y \implies T = T^{**}$$
    Por lo tanto, podemos finalizar la demostración de $ii)$ con $\|T^{**}\| \le \|T^*\| \iff \|T\| \le \|T^*\|$ y en consecuencia $\|T\| = \|T^*\|$.
\end{proof}
\end{prop}

\begin{prop}
    Sean $S, T \in BL(\mathcal{H})$, entonces se tiene:
    \begin{itemize}
        \item $(S + T)^* = S^* + T^*$
        \item $(\alpha T)^* = \bar{\alpha} T^*$
        \item $(S \circ T)^* = T^* \circ S^*$
        \item $\|T^* \circ T\| = \|T \circ T^*\| = \|T\|^2$ \\
    \end{itemize}
\end{prop}

\begin{ejercicio}
    Demostrar 
    \begin{itemize}
        \item[i)] $N(T^*) = R(T)^\perp$
        \item[ii)] $N(T) = R(T^*)^\perp$ \\
    \end{itemize}
    \noindent
    Primero tengamos claro que estos conjuntos son el rango y el núcleo, es decir
    \begin{equation*}
        N(T) = \{x \in \mathcal{H} \mid T(x) = 0\}\quad \text{y} \quad  R(T) = \{y \in \mathcal{H} \mid y = T(x) \text{ para algún } x \in \mathcal{H}\}
    \end{equation*}
    Para demostrar $ii)$, supongamos  $i)$, es decir, $N(T^*) = R(T)^\perp$, entonces 
    $$N(T^{**}) = R(T^*)^\perp$$
    Como $T^{**} = T$, tenemos finalmente $N(T) = R(T^*)^\perp$. \\

    \noindent
    Demostremos ahora $i)$, para ello veámos los conjutnos
    \begin{align*}
        N(T^*) &= \{x \in \mathcal{H} \mid T^*(x) = 0\} \\
        R(T)^\perp &= \{x \in \mathcal{H} \mid (x, a)=0 \ \forall a \in R(T)\}= \{x \in \mathcal{H} \mid (x, T(y)) = 0 \quad \forall y \in \mathcal{H}\}
    \end{align*}
    entonces
    $$x \in N(T^*) \implies T^*(x) = 0 \implies (x, T(y)) = (T^*(x), y) = (0,y)=0 \quad \forall y \in \mathcal{H} \implies x \in R(T)^\perp$$
    donde se tiene $N(T^*) \subseteq R(T)^\perp$. Veámos el viceverso, por lo que si $x \in R(T)^\perp$ se cumple que
    $$(x, T(y)) = 0 \quad \forall y \in \mathcal{H} \implies (T^*(x), y) = 0 \quad \forall y \in \mathcal{H} \implies T^*(x) = 0$$
\end{ejercicio}

\begin{definicion}
    Sea $T \in BL(\mathcal{H})$ se dice \textbf{autoadjunto} si $T = T^*$, es decir
    \[ (x, T(y)) = (T(x), y) \quad \forall x, y \in \mathcal{H} \]
\end{definicion}

\begin{ejemplo}
    Sea $\ell_2$ con el producto interno usual, es decir, $(x, y) = \sum\limits_{i=0}^{+\infty} x_i \cdot y_i$, y $\lambda = \{\lambda_i\}_{i=1}^\infty \in \ell_\infty$. Definamos 
    \Func{T}{\ell_2}{\ell_2}{x}{(\lambda_1 x_1, \lambda_2 x_2, \dots, \lambda_m x_m, \dots)}
    y queremos ver que es autoadjunto. Para ello, veámos primero $T\in BL(\ell_2)$, para lo cual se tiene
    \begin{itemize}
        \item \tbf{Linealidad:} $T(\alpha x + \beta y) = (\lambda_1 (\alpha x_1 + \beta y_1), \lambda_2 (\alpha x_2 + \beta y_2), \dots) = \alpha T(x) + \beta T(y)$
        \item \tbf{Continuidad:} tenemos que 
        \begin{equation*}
            \|T(x)\|^2 = (T(x), T(x)) = \sum_i \lambda_i^2 x_i^2 \le \sup_i \lambda_i^2 \cdot \sum_i |x_i|^2 = \|\lambda\|_{\ell_\infty}^2 \cdot \|x\|_{\ell_2}^2
        \end{equation*}
    \end{itemize}
    Como $\lambda$ es acotada: sea $\lambda_k : |\lambda_k| \ge \|\lambda\|_{\ell_\infty} - \varepsilon$, con $\varepsilon > 0$. Entonces, para $\tilde{x} = (\delta_{ki})_i$ (vector de la base canónica), se tiene:
    \[ \|T(\tilde{x})\| = \sqrt{0^2 + 0^2 + \dots + |\lambda_k|^2 + \dots} = |\lambda_k|= \ge \|\lambda\|_{\ell_\infty} - \varepsilon \xrightarrow{\varepsilon \to 0} \|\lambda\|_{\ell_\infty} \implies \|T\| = \|\lambda\|_{\ell_\infty} \]
    Tenemos ya que $T$ es un operador lineal y acotado, veamos ahora que es autoadjunto.
    \[ (x, T(y)) = \sum_i x_i \lambda_i y_i = \sum_i (x_i \lambda_i) y_i = (T(x), y) \quad \forall x, y \in \mathcal{H} \]

\end{ejemplo}


\begin{ejercicio}
    Sea $\mathcal{H}$ un espacio de Hilbert y sea $M \subseteq \mathcal{H}$ un subespacio cerrado y $P_M: \mathcal{H} \to \mathcal{H}$ la proyección. Entonces:
    \begin{itemize}
        \item[i)] $P_M$ es autoadjunto.
        \item[ii)] $(P_M(x), x) \ge 0$.
        \item[iii)] $\|P_M\| = 1$ (si $M \neq \{0\}$).
    \end{itemize}
    \begin{description}
        \item[$i)$] Para demostrar que $P_M$ es autoadjunto, es decir, que $(x, P_M(y)) = (P_M(x), y) \quad \forall x, y \in \mathcal{H}$, usaremos la descomposición ortogonal de $\mathcal{H}$ respecto a $M$ y $M^\perp$. Dado que $M$ es un subespacio cerrado, entonces $\mathcal{H} = M \oplus M^\perp$, por lo tanto $\forall x, y \in \mathcal{H}$
        \begin{equation*}
            x = P_M(x) + P_{M^\perp}(x) \quad \text{y} \quad y = P_M(y) + P_{M^\perp}(y)
        \end{equation*}
        entonces
        \begin{align*}
            (x, P_M(y)) &= (P_M(x) + \underbrace{P_{M^\perp}(x)}_{\in M^\perp}, \underbrace{P_M(y)}_{\in M}) = (P_M(x), P_M(y))+(P_{M^\perp}(x), P_M(y))= (P_M(x), P_M(y))\\
            (P_M(x), y) &= (P_M(x), P_M(y) + P_{M^\perp}(y)) = (P_M(x), P_M(y)) 
        \end{align*}
        donde hemos usado que $P_M(y) \in M$ y $P_{M^\perp}(x) \in M^\perp$ son ortogonales, por lo que $(P_{M^\perp}(x), P_M(y)) = 0$. De lo cual $(x, P_M(y)) = (P_M(x), y)$.
        \item[$ii)$] Para demostrar que $(P_M(x), x) \ge 0$, usaremos la misma descomposición ortogonal que en el apartado anterior, es decir, $x = P_M(x) + P_{M^\perp}(x)$, entonces
        \begin{equation*}
            (P_M(x), x) = (P_M(x), P_M(x) + P_{M^\perp}(x)) = (P_M(x), P_M(x)) = \|P_M(x)\|^2 \ge 0
        \end{equation*}
        \item[$iii)$] Para demostrar que $\|P_M\| = 1$ (si $M \neq \{0\}$), usaremos la misma descomposición ortogonal que en los apartados anteriores, es decir, $x = P_M(x) + P_{M^\perp}(x)$, entonces
        \begin{equation*}
            \|P_M(x)\| \overset{(*)}{=} \|x\| - \underbrace{\|P_{M^\perp}(x)\|}_{\ge 0} \le \|x\| \implies \|P_M\| \le 1
        \end{equation*}
        donde en $(*)$ hemos usado el teorema de pitágoras, que indica que si $M$ y $M^\perp$ son ortogonales, entonces $\|x\|^2 = \|P_M(x)\|^2 + \|P_{M^\perp}(x)\|^2$. \\
        \noindent
        Sea $x \in M = \overline{M}$, entonces $P_M(x) = x \implies \|P_M\| = \sup\limits_{\|x\|\neq 0} \frac{\|P_M(x)\|}{\|x\|} =1$, para algún $x \in M$ con $\|x\|\neq 0$, que lo hay por hipótesis de que $M \neq \{0\}$, por lo que $\|P_M\| = 1$.
    \end{description}
\end{ejercicio}

\begin{teo}
    Sea $T \in BL(\mathcal{H})$ autoadjunto, entonces 
    \begin{equation*}
        \|T\| = \sup_{x \neq 0} \frac{|(T(x), x)|}{\|x\|^2} = \sup_{\|x\|=1} |(T(x), x)|
    \end{equation*}
\begin{proof}
    Antes de demostrar la igualdad siguiente
    \begin{equation*}
        \|T\|  = \sup_{\|x\|=1} |(T(x), x)|
    \end{equation*}
    Vamos a recordar porque la igualdad 
    $$\sup\limits_{x \neq 0} \dfrac{|(T(x), x)|}{\|x\|^2} = \sup\limits_{\|x\|=1} |(T(x), x)|$$ 
    es obvia. Para ello, sea $x \neq 0$, podemos definir un vector normalizado $u = \frac{x}{\|x\|}$, de tal manera que $\|u\| = 1$. Si sustituimos esto en la expresión del lado izquierdo:  
    \begin{itemize}
        \item Por la linealidad del operador $T$, tenemos que $T(x) = T(\|x\| \cdot u) = \|x\| T(u)$.
        \item Por las propiedades del producto escalar, se tiene que
        \begin{equation*}
            (T(x), x) = (\|x\| T(u), \|x\| u) = \|x\|^2 (T(u), u)
        \end{equation*}
        \item Sustituyendo en el cociente, se obtiene
        \begin{equation*}            
            \frac{|(T(x), x)|}{\|x\|^2} = \frac{\|x\|^2 |(T(u), u)|}{\|x\|^2} = |(T(u), u)| \\
        \end{equation*}
    \end{itemize}

    \noindent
    Ahora centrémonos en la igualdad complicada. Sea $x \in \mathcal{H}$ tal que $\|x\| = 1$, entonces \\ \mbox{$|(Tx, x)| \le \|Tx\| \cdot \|x\| \le \|T\| \cdot \|x\|^2 = \|T\|$}, por lo tanto
    \[ \sup_{\|x\|=1} |(T(x), x)| \le \|T\| \]
    Mostremos ahora la desigualdad inversa. Sea 
    $$M = \sup_{\|x\|=1} |(T(x), x)| = \sup_{x \neq 0} \frac{|(T(x), x)|}{\|x\|^2}$$
    tenemos que $|(T(x), x)| \le M \|x\|^2 \quad \forall x \in \mathcal{H}$. Haciendo los siguientes desarrollos, se obtiene
    \begin{align*}
        (T(x+y), x+y) &= (Tx + Ty, x+y) = (Tx, x) + (Ty, y) + (Tx, y) + (x, Ty)= \\
        &\overset{(*)}{=} (Tx, x) + (Ty, y) + 2(Tx, y) \\
        (T(x-y), x-y) &= (Tx, x) + (Ty, y) - 2(Tx, y)
    \end{align*}
    donde en $(*)$ hemos usado que $T$ es autoadjunto, y restando las dos ecuaciones anteriores, se obtiene
    \begin{equation*}
        (T(x+y), x+y) - (T(x-y), x-y) = 4(Tx, y)
    \end{equation*}
    Ahora aplicanddo los desarrollos anteriores y la definición de $M$, tenemos que
    \begin{align*}
        (Tx, y) &= \frac{1}{4} \left( (T(x+y), x+y) - (T(x-y), x-y) \right)\leq \\
        &\le \frac{1}{4} \left( |(T(x+y), x+y)| + |(T(x-y), x-y)| \right) \leq\\
        &\le \frac{1}{4} M \left( \|x+y\|^2 + \|x-y\|^2 \right) \overset{(*)}{=} \frac{1}{2} M (\|x\|^2 + \|y\|^2) \quad \forall x, y \in \mathcal{H}
    \end{align*}
    donde en $(*)$ hemos usado la ley del paralelogramo, en Teorema \ref{teo:paralelogramo}. Sea $x : \|x\|=1$ y $y = \dfrac{Tx}{\|Tx\|}$ entonces 
    $$(Tx, y) = \left( Tx, \frac{Tx}{\|Tx\|} \right) \le \frac{1}{2} M \cdot 2$$
    de lo cual
    \begin{equation*}
        (T(x), y) = \dfrac{(T(x),T(x))}{\|T(x)\|}=\frac{\|Tx\|^2}{\|Tx\|} = \|Tx\| \le M \quad \forall x : \|x\|=1
    \end{equation*}
    entonces aplicando la definición de norma de operador, tenemos lo buscado
    \begin{equation*}
        \|T\| = \sup_{\|x\|=1} \|Tx\| \le M
    \end{equation*}
\end{proof}
\end{teo}

\begin{prop}
    Sea $T\in BL(\mathcal{H})$, entonces 
    \begin{equation*}
        T \text{ biyectivo} \iff T^* \text{ biyectivo}
    \end{equation*} 
\end{prop}

\section{Operadores compactos}

\begin{definicion}
    Sea $X$ un espacio normado, un subconjunto $K \subseteq X$ es \textbf{compacto} si toda sucesión $\{x_n\}$ en $K$ admite una subsucesión convergente a un punto de $K$.
\end{definicion}

\begin{definicion}
    Sea $X$ un espacio normado, un subconjunto $A \subseteq X$ es \textbf{relativamente compacto} si la clausura de $A$ es compacta.
\end{definicion}

\begin{definicion}
    Sean $X, Y$ espacios de Banach y sea $T\in BL(X,Y)$. $T$ se dice \textbf{operador compacto} si $T(C)$ es \underline{relativamente compacto} para cada conjunto $C$ acotado. \\
    \noindent
    Análogamente, $T$ es \tbf{operador compacto} si para cada $\{x_n\}$ sucesión acotada, $T(x_n)$ admite una subsucesión convergente.
\end{definicion}


\begin{ejemplo}
    \begin{enumerate}
        \item Si $T: X \to Y$ es lineal y $X, Y$ son espacios normados de dimensión finita, entonces $T$ es compacto, veámoslo.
        
        $X, Y$ son espacios de dimensión finita $\implies T$ es continua 
        $\equiv$ acotada $\implies$ manda sucesiones acotadas en sucesiones acotadas en $Y \simeq \mathbb{R}^n \implies$ admite subsucesiones convergentes.
        
        En general, si $X$ un espacio de Banach e $Y$ un espacio normado de \underline{dimensión finita}, si $T \in BL(X, Y)$, entonces $T$ es un operador compacto.

        \item $Id: X \to X$ nunca es compacto para cualquier $X$ espacio normado de dimensión infinita. De hecho: la bola unitaria no es compacta.
    \end{enumerate}
\end{ejemplo}

\observacion{
    $\mathcal{K}(X, Y) := \{ T \in BL(X, Y) \text{ compacto} \} \subseteq BL(X, Y)$ es un subespacio de $BL(X, Y)$. \\
}

\begin{teo}
    Sea $\{T_n\} \subseteq \mathcal{K}(X, Y)$ una sucesión de operadores compactos tales que \mbox{$T_n \to T\in BL(X, Y)$}, entonces $T$ es compacto.
\begin{proof}
    Sea $(x_1, \dots, x_m, \dots)$ una sucesión acotada en $X$, mostremos que $T(x_k)$ admite una subsucesión convergente.

    $T_1$ es compacto $\implies (T_1(x_k))_k$ tiene una subsucesión convergente, llamémosla $(T_1(x_k^1))$.

    $T_2$ es compacto $\implies (T_2(x_k^1))$ tiene una subsucesión convergente $\implies (T_2(x_k^2))$ es convergente.

    Iterando: sea $(x_k^n)$ una subsucesión de $(x_k^{n-1}) \implies T_j(x_k^n)$ es convergente para cada $j \le n$.

    Sea $\{x_k^k\}$ la \textbf{sucesión diagonal} $\implies \{T_j(x_k^k)\}$ es convergente $\forall j$. Mostremos que $\{T(x_k^k)\}$ es convergente (o de Cauchy):

    \[ \|T(x_k^k) - T(x_m^m)\| \le \|T(x_k^k) - T_j(x_k^k)\| + \|T_j(x_k^k) - T_j(x_m^m)\| + \|T_j(x_m^m) - T(x_m^m)\| \le \]
    \[ \le \underbrace{\|T - T_j\|}_{\le \varepsilon} \cdot \|x_k^k\| + \underbrace{\|T_j(x_k^k) - T_j(x_m^m)\|}_{\le \varepsilon} + \underbrace{\|T_j - T\|}_{\le \varepsilon} \cdot \|x_m^m\| \le 3\varepsilon \to 0 \]
\end{proof}
\end{teo}

\observacion{
    Si $(T_n)$ es tal que $R(T_n)$ es de dimensión finita y $T_n \to T$, entonces $T$ es compacto.
}


\begin{prop}
    Sean $X$ un espacio de Banach e $Y$ un espacio de Hilbert, y sea $T \in K(X, Y)$. Entonces $T$ es el límite de $(T_n)_n$, donde $(T_n)_n$ es una sucesión de operadores con $R(T_n)$ de dimensión finita.
\end{prop}

\begin{ejercicio}
    Dada la función
    \Func{A}{L^2((0, 1))}{L^2((0, 1))}{u(t)}{Au(t):=a(t) \cdot u(t)}
    para todo $t \in (0, 1)$, con $a: (0, 1) \longrightarrow \mathbb{R}$ continua y $a \neq 0$. Demuestra que $A$ es un operador lineal y continuo, pero no compacto. \\

    \noindent
    Demostraremos la no compacidad, la definición nos dice que debemos encontrar una sucesión acotada de funciones $(e_n)$ tal que, al aplicarles el operador, la sucesión de imágenes $(A(e_n))$ no tenga ninguna subsucesión convergente. \\

    \noindent
    Como $a$ es no nula, entonces $\exists t_0 \in (0, 1) : a(t_0) \neq 0$, y como $a$ es continua tenemos un intervalo dado por 
    $$\exists\delta>0 \text{ tal que } |a(t)| \ge \frac{1}{2} |a(t_0)| > 0 \quad \forall t \in (t_0 - \delta, t_0 + \delta)$$
    Ahora usaremos $(e_n)$ una base ortonormal (o una sucesión ortonormal) que está concentrada o vive dentro de ese pequeño intervalo $(t_0 - \delta, t_0 + \delta)$. Como es una base ortonormal, se sabe por propiedades de los espacios de Hilbert que
    \begin{itemize}
        \item $(e_n)$ es una sucesión acotada, de hecho $\|e_n\| = 1$ para cada $n$.
        \item $(e_n)$ no tiene ninguna subsucesión convergente, de hecho, la distancia entre dos elementos distintos es exactamente $\sqrt{2}$, por lo que no pueden acercarse entre sí, veámoslo
        \begin{equation*}
            \|e_n - e_m\|^2 = (e_n, e_n) - (e_n, e_m) - (e_m, e_n) + (e_m, e_m)= \|e_n\|^2 -0-0+ \|e_m\|^2 = 2 \quad \forall n \neq m
        \end{equation*}
    \end{itemize}
    entonces tenemos 
    \begin{align*}
        \|A(e_n) - A(e_m)\|_{L^2}^2 &= \int_{0}^{1} |a(t) \cdot (e_n(t) - e_m(t))|^2 dt \geq \frac{1}{4} |a(t_0)|^2 \int_{t_0 - \delta}^{t_0 + \delta} |e_n(t) - e_m(t)|^2 dt \geq \\
        &\overset{(*)}{\geq} \frac{1}{4} |a(t_0)|^2 \cdot 2 \geq \frac{1}{2} |a(t_0)| > 0 \quad \forall n, m
    \end{align*}
    donde únicamente hemos aplicado lo deducido anteriormente y la norma de $L^2$, y en $(*)$ hemos usado este razonamiento
    \begin{align*}
        \int_{t_0 - \delta}^{t_0 + \delta} |e_n(t) - e_m(t)|^2 dt &=\int_{t_0 - \delta}^{t_0 + \delta} \left( |e_n(t)|^2 - 2e_n(t)e_m(t) + |e_m(t)|^2 \right) dt = \\
        &=\int_{t_0 - \delta}^{t_0 + \delta} |e_n(t)|^2 dt - 2\int_{t_0 - \delta}^{t_0 + \delta} e_n(t)e_m(t) dt+ \int_{t_0 - \delta}^{t_0 + \delta} |e_m(t)|^2 dt= \\
        &=1-2\cdot 0+1=2
    \end{align*}
    donde se ha usado que 
    \begin{align*}
        \int_{t_0 - \delta}^{t_0 + \delta} |e_n(t)|^2 dt &= \|e_n\|_{L^2}^2 = 1 \\
        \int_{t_0 - \delta}^{t_0 + \delta} |e_m(t)|^2 dt &= \|e_m\|_{L^2}^2 = 1 \\
        \int_{t_0 - \delta}^{t_0 + \delta} e_n(t)e_m(t) dt &= (e_n, e_m)_{L^2} = 0
    \end{align*}
    lo que hemos demostrado es que la sucesión de esta base ortonormal, no tiene subsucesiones convergentes, aun cuando $(e_n)$ está acotada. Por lo tanto, $A$ no es compacto.
\end{ejercicio}

\begin{ejercicio}
    Sea $T: \mathcal{H} \longrightarrow \mathcal{H}$ lineal, entonces
    \begin{equation*}
        T \text{ es compacto } \iff T^* \text{ es compacto}
    \end{equation*}
\end{ejercicio}



